% !TeX root = ../main.tex

\chapter{Cascade Routing between Direct and Pipeline Flows}
\label{ch:04}

The pipeline of Chapter~\ref{ch:03} is not uniformly better than simply
asking the model for the module. On a substantial class of tasks it is
\emph{worse}: its planning and reuse scaffolding invents structure that a
self-contained module does not need, and derails a generation that a bare
prompt would have gotten right. This chapter presents the second
contribution of the thesis: a per-task \emph{cascade router} that predicts,
before spending pipeline-scale compute, which of two flows a task should
take --- the full agentic pipeline, or a cheap \emph{direct} single-shot
generation. Section~\ref{sec:router-motivation} establishes why the two
flows are complementary; Section~\ref{sec:router-classes} defines the task
classes and the oracle labels used to evaluate routing;
Section~\ref{sec:router-design} describes the two-tier cascade; and
Section~\ref{sec:router-metrics} defines the routing-specific metrics used
in Section~\ref{sec:results-routing}.

\section{Motivation: Two Complementary Flows}
\label{sec:router-motivation}

Running \emph{both} flows on all 60 tasks (the \textsc{Hybrid} arm of
Table~\ref{tab:run-inventory}) reveals a clean division of labor. With ten
samples per flow, the pipeline alone solves 20 tasks and the direct flow
alone solves 22 --- but their union solves 30, because the two sets overlap
only partially. Wrapper and integration modules pass under the pipeline,
which supplies the exact interfaces of the submodules they must instantiate
and wire; self-contained algorithmic modules pass under direct generation,
which lets the model write the datapath in one coherent piece without the
pipeline's decomposition machinery second-guessing it. Neither flow
dominates: each one's strength is the other's failure mode.

The obvious exploitation of complementarity --- run both flows on every
task --- is also the most expensive one, and it wastes most of its budget:
every direct sample spent on a wrapper module and every pipeline run spent
on a self-contained ALU is compute that predictably converts to nothing.
A router that sends each task to its matching flow can, in principle, reach
the union coverage at roughly half the cost. The routing problem is to make
that prediction \emph{before} the expensive flow has been paid for.

\section{Task Classes and Oracle Labels}
\label{sec:router-classes}

We formalize the division as two task classes:

\begin{description}
  \item[Class A (structural).] Wrapper, integration, and memory-encapsulation
    modules: the module's own logic is thin; its difficulty is wiring
    externally defined submodule interfaces correctly. Class~A tasks want
    the \textbf{pipeline}.
  \item[Class B (algorithmic).] Self-contained datapath or control modules
    --- cipher rounds, CRC, instruction decode, address generation: the
    difficulty is the logic itself, not integration. Class~B tasks want
    \textbf{direct} generation.
\end{description}

For evaluation only, each task receives an oracle label computed from its
\emph{golden} implementation: the reference RTL is synthesized with Yosys,
submodule instances are subtracted, and the task is labeled~A if its own
synthesized logic falls below a 250-cell threshold (or it is a recognized
regular-array wrapper), otherwise~B. The label also fixes the task's
\emph{target flow} (A$\to$pipeline, B$\to$direct). Oracle labels use golden
RTL that no generation-time component may see; they exist so that routing
accuracy can be scored, and they define the ceiling arm (\texttt{oracle}) in
the ablation of Section~\ref{sec:results-routing}.

\section{The Two-Tier Cascade}
\label{sec:router-design}

\begin{figure}[htbp]
  \centering
  \begin{tikzpicture}[
      font=\small,
      node distance=9mm and 14mm,
      box/.style={draw, rounded corners=2pt, align=center, inner sep=5pt,
                  minimum height=8mm},
      dec/.style={draw, diamond, aspect=2.2, align=center, inner sep=1.5pt},
      flow/.style={draw, rounded corners=2pt, align=center, inner sep=5pt,
                   minimum height=8mm, very thick},
      arr/.style={-{Stealth[length=2.5mm]}},
    ]
    \node[box] (spec) {task specification};
    \node[box, below=of spec] (feat) {Tier-0: feature extraction\\(LLM judge over the spec)};
    \node[dec, below=of feat] (pre) {pre-route\\decision};
    \node[flow, below left=13mm and 16mm of pre] (direct) {\textbf{direct flow}\\single-shot generation\\+ repair loops};
    \node[flow, below right=13mm and 16mm of pre] (pipe) {\textbf{pipeline flow}\\plan $\to$ generate\\+ repair loops};
    \node[box, below=17mm of pre] (plan) {Tier-1: run planner once,\\probe the structured plan};
    \node[dec, below=of plan] (probe) {plan\\verdict};
    \draw[arr] (spec) -- (feat);
    \draw[arr] (feat) -- (pre);
    \draw[arr] (pre) -| node[above, pos=0.25]{algorithmic (B)} (direct);
    \draw[arr] (pre) -| node[above, pos=0.25]{structural (A)} (pipe);
    \draw[arr] (pre) -- node[right]{uncertain} (plan);
    \draw[arr] (plan) -- (probe);
    \draw[arr] (probe) -| node[above, pos=0.25]{self-contained} (direct);
    \draw[arr] (probe) -| node[above, pos=0.25]{delegates (keep plan)} (pipe);
  \end{tikzpicture}
  \caption[Cascade router decision flow]{The two-tier cascade router. Tier-0
    routes confident tasks from spec-level features alone; only uncertain
    tasks pay for a planner call, whose structured plan Tier-1 probes. A
    task routed to direct after planning discards the plan (a ``wasted
    plan'').}
  \label{fig:router-flow}
\end{figure}

The router is a cascade of two increasingly expensive tests
(Figure~\ref{fig:router-flow}); a task exits at the first tier that is
confident.

\paragraph{Tier-0: spec pre-route.}
A feature extractor reads the specification and emits seven structured
judgments: whether the module delegates to submodules, is a memory
encapsulation, or is thin pass-through control (all class-A signals);
whether it contains a finite-state machine, clock-domain crossing, or a
datapath algorithm (class-B signals); and an estimate of the state bits the
module itself must hold, plus a confidence value. The production extractor
is a small LLM call that returns this JSON verdict (a regex-based keyword
extractor over the same features serves as the ablation baseline ---
Section~\ref{sec:results-routing}); the decision rule is deliberately
transparent: any algorithmic signal routes to \emph{direct}, else any
structural signal routes to \emph{pipeline}, else --- or when the
extractor's confidence is below threshold --- the task falls through as
\emph{uncertain}. One subtlety the LLM extractor handles and keywords
cannot: a module can \emph{sound} algorithmic (``preliminary decoding'')
while actually being a wrapper that instantiates the real decoder, and only
whole-spec reading resolves that.

\paragraph{Tier-1: plan probe.}
Uncertain tasks run the planner once --- the cheapest informative pipeline
component --- and the router probes the resulting structured plan's own
natural-language requirement summary. Decisive delegation phrasing
(``wraps'', ``encapsulates'', ``reuses existing'', or instantiation language
without any own-implementation verb) keeps the task in the pipeline, plan
in hand; clear self-contained-algorithm phrasing routes it to direct,
\emph{discarding} the plan. The asymmetry of the default is intentional:
the plan is already a sunk cost, so the probe bails to direct only when the
module is clearly self-contained. A plan generated and then discarded is
counted as a \emph{wasted plan} --- the cascade's overhead metric.

\paragraph{The direct flow.}
The direct arm is not a bare completion: it shares the generator's repair
machinery. A directly generated candidate that fails the internal lint
enters the same syntax repair loop, and (in arms where functional repair is
enabled) the same functional repair loop, with the plan-dependent prompt
sections simply absent. This matters for the routing economics: direct
generation costs roughly 10--12\,k tokens per sample against the pipeline's
90--153\,k, so every correctly-routed class-B task buys nearly an order of
magnitude of compute reduction at equal or better quality.

\section{Routing Metrics}
\label{sec:router-metrics}

Section~\ref{sec:results-routing} evaluates routing arms on four axes:
\emph{coverage} (solved tasks and pass@k, against the \textsc{Hybrid} union
as the attainable ceiling); \emph{routing accuracy} (the confusion between
oracle class and chosen flow, where both misroute directions ---
A$\to$direct and B$\to$pipeline --- forfeit a would-be pass);
\emph{overhead} (wasted plans); and \emph{efficiency} (estimated tokens and
summed compute per solved task). The design goal is stated once and
measured throughout: reach the union coverage of both flows at the lowest
compute, with misroutes near zero.
